\section{Ảnh thành token}
\label{sec:ch11-anh-thanh-token}

\index{Vision Transformer!ảnh thành token}%
Một ảnh không tự nhiên là một câu. Pixel đứng cạnh nhau có quan hệ mà token
trong câu không có, và chương~\ref{ch:bias-quy-nap} vừa dành cả một chương để
đếm giá của quan hệ ấy. ViT không phủ nhận điều đó. Nó chọn một đường khác:
gói nhiều pixel cạnh nhau thành một \tn{mảnh ảnh}{patch}, rồi để self-attention
học quan hệ giữa các gói. Khác với \tn{tích chập}{convolution}, nó không buộc
\tn{tính cục bộ}{locality} giữa các gói, cũng không có \tn{chia sẻ trọng
số}{weight sharing} theo từng lệch pixel.

Gọi ảnh đầu vào là \(\mat{X} \in \R^{H \times W \times C}\), patch vuông có
cạnh \(P\). Bỏ qua biên cho gọn, ảnh tách thành

\begin{equation}
\label{eq:ch11-so-patch}
N = \frac{HW}{P^2}
\end{equation}

mảnh. Mỗi mảnh có \(P^2C\) số, được làm phẳng rồi nhân với cùng ma trận
\(\mat{E}\) để thành một vector \(D\) chiều. Đây là phép chiếu tuyến tính, không
phải tích chập: hai patch ở hai chỗ khác nhau đi qua cùng \(\mat{E}\), nhưng
\(\mat{E}\) nhìn trọn patch 8x8 một lần chứ không trượt từng pixel trên ảnh.

Hình~\ref{fig:ch11-anh-thanh-token} vẽ đúng phiên bản tôi chạy sau này. CIFAR-10
là ảnh RGB 32x32. Patch 8x8 tạo 16 patch, mỗi patch có 192 số. Một
\tn{token lớp}{class token} đứng trước 16 token ấy. Con số patch nhỏ hơn 16x16
của nhan đề bài báo, vì nếu giữ 16x16 trên ảnh 32x32 thì cả ảnh chỉ còn bốn
token và câu hỏi về attention gần như biến mất.

% Width measured with \savebox: 302.15pt against this book's 441.02pt measure,
% leaving 138.87pt. Do not estimate it from the page (quyet dinh 67).
\begin{figure}[htbp]
  \centering
  % Ảnh 32x32 ở trên được chia thành lưới 4x4 patch 8x8. Mỗi patch được làm
% phẳng rồi qua cùng một phép chiếu tuyến tính. Dưới ảnh là dãy 17 token:
% class token đứng trước 16 token ảnh. Positional embedding được cộng vào mọi
% ô trong dãy. Hình chỉ vẽ input của encoder, không vẽ attention.
% Mono-safe: patch khác nhau bằng độ xám xen kẽ; cấu trúc nằm ở lưới và mũi tên.
\begin{tikzpicture}[
  >=Stealth,
  anh/.style={draw, minimum size=6mm, inner sep=0pt},
  patchdam/.style={anh, fill=black!18},
  patchnhat/.style={anh, fill=black!4},
  token/.style={draw, rounded corners=1.5pt, minimum width=5.6mm,
                minimum height=5mm, font=\tiny, fill=white},
  cls/.style={token, fill=black!16, font=\tiny\bfseries},
  nhan/.style={font=\footnotesize},
  nho/.style={font=\scriptsize, black!65},
]
  \node[nhan] at (1.5, 3.45) {Ảnh RGB 32x32};
  \foreach \r in {0,1,2,3} {
    \foreach \c in {0,1,2,3} {
      \pgfmathtruncatemacro\k{4*\r+\c+1}
      \ifodd\k
        \node[patchdam] (p\k) at (\c, {2.45-\r}) {};
      \else
        \node[patchnhat] (p\k) at (\c, {2.45-\r}) {};
      \fi
    }
  }
  \draw[black!55] (-0.38,-0.93) rectangle (3.38,2.83);
  \node[nho, align=center] at (1.5,-1.48) {16 mảnh ảnh 8x8\\mỗi mảnh có 192 số};

  \draw[-{Stealth[length=2mm]}, black!65] (1.5,-1.95) -- (1.5,-2.7);
  \node[nho, anchor=west] at (1.75,-2.32)
    {làm phẳng, chiếu bằng cùng ma trận \(\mat{E}\)};

  \node[nhan] at (5.0,-3.2) {Chuỗi vào encoder};
  \node[cls] (c) at (-0.05,-4.05) {class};
  \foreach \k in {1,...,16} {
    \node[token] (t\k) at ({0.42+0.58*\k},-4.05) {\(p_\k\)};
  }
  \draw[black!45, rounded corners=2pt] (-0.48,-4.5) rectangle (9.85,-3.6);
  \node[nho, align=center] at (4.7,-5.0)
    {17 token: một token lớp và 16 token ảnh.\\
     Mỗi token nhận thêm một embedding vị trí học được.};
\end{tikzpicture}

  \caption{Ảnh không đi thẳng vào encoder. Nó được chia thành 16 mảnh ảnh,
  mỗi mảnh được làm phẳng và chiếu bằng cùng một ma trận, rồi thành 16 token.
  Token lớp đứng đầu dãy, nên encoder nhận 17 token. Cái hình cố ý không vẽ
  một cửa sổ cục bộ nào: sau đường cắt patch, mọi token có thể nhìn mọi token.}
  \label{fig:ch11-anh-thanh-token}
\end{figure}

Chữ \enquote{cùng ma trận} dễ làm người ta gọi đây là chia sẻ trọng số rồi đi
tiếp. Không đúng theo nghĩa chương trước cần. Một kernel 3x3 được dùng lại ở
mỗi lệch một pixel, nên nó buộc một con mèo ở góc và một con mèo ở giữa sinh ra
cùng kiểu đặc trưng. Phép chiếu patch chỉ buộc bốn mươi tám hay một trăm chín
mươi hai pixel \emph{bên trong một patch} dùng cùng phép chiếu. Nó không nói
patch thứ nhất và patch thứ mười sáu phải có quan hệ nào.

Điều đó là điểm xuất phát của ViT, không phải một khuyết điểm nhỏ cần giấu.
Mục sau nói cách mô hình biết patch nào đang ở đâu.
